\section{Relation to neighboring frameworks}
\label{sec:related}

The proposal intersects several mature literatures.  Its identity depends on
what it does not conflate with proof-bearing recovery.

\subsection{Structural realism and theory change}

Epistemic structural realism seeks a stable commitment across theory change,
while ontic versions make stronger claims about the priority or exclusivity of
structure.  Worrall's account of structural continuity supplies a central
motivation \cite{Worrall1989}.  Van Fraassen and Roberts expose problems of
content, triviality, and hierarchy when ``structure'' is insufficiently typed
\cite{VanFraassen2006,Roberts2011}.  Our profiles are not a criterion for
fundamentality.  They identify the exact claim, observable class, and theorem
that survive one transformation.  This gives structural continuity an
inferential test without converting it into a complete ontology.

\subsection{Gauge symmetry and surplus structure}

Philosophical treatments of gauge theory distinguish transformations that
relate physically identical presentations from transformations with a
different interpretive role \cite{Belot2003}.  Sophistication treats
symmetries as isomorphisms rather than necessarily reducing to invariant
objects \cite{Dewar2019}.  Nguyen, Teh, and Wells use category-theoretic
comparisons to show that surplus gauge structure can be indispensable for a
rich collection of local fields \cite{NguyenTehWells2020}.  This is the
closest precedent for our claim that one property can be preserved while
another becomes unavailable.  We extend the diagnostic from gauge reduction
to indexed regulator, RG, continuum, reconstruction, and spectral arrows, and
attach quantitative moduli and external recovery theorems.

\subsection{Equivalence and duality}

Formal criteria alone may not settle theoretical equivalence because
interpretation also matters \cite{Coffey2014}; algebraic weak equivalence can
likewise fail to establish physical equivalence in QFT
\cite{Lupher2018}.  De Haro and Butterfield distinguish bare theories, model
roots, specific structure, and interpretation, and represent exact duality as
an isomorphism of model triples \cite{DeHaroButterfield2018}.  Our exact
subcalculus is compatible with that schema.  Most arrows of interest here are
not dualities: a cutoff removal is indexed and singular, a coarse graining is
not invertible, and OS reconstruction is theorem-backed rather than a mere
redescription.  Calling all of them equivalences would erase the central
problem.

\subsection{Perspectivism, models, and perspicuity}

Perspectival and model-based accounts emphasize that scientific access is
conditioned by instruments, models, and aims.  North's recent account of
perspicuous representation argues that some physical representations reveal
underlying structure better than others and that this is not wholly a matter
of stipulation \cite{North2026}.  Proof-bearing dynamics asks a complementary
question.  A representation may be perspicuous for discovering an estimate
yet insufficient for transporting it.  Conversely, an imperspicuous
intermediate formalism may be acceptable when a recovery theorem closes the
relevant property channel.  Perspicuity concerns what a formulation makes
visible; certificates concern what an inference establishes.

\subsection{RG, EFT, emergence, and limits}

Wilson and Kogut's account makes scale transformation a technical operation
on theory space \cite{WilsonKogut1974}.  Philosophical work has used RG and
EFT to analyze emergence, reduction, autonomy of scales, and realism under
infinite idealization \cite{Butterfield2011,Rivat2021}.  Koberinski and Fraser
show that different RG variants play different epistemic roles in particle
physics and condensed matter, and caution against identifying RG generically
with coarse graining \cite{KoberinskiFraser2023}.  Our taxonomy adopts that
domain specificity.  It contributes a proof-level question: which observable,
positivity, locality, and limit certificates are actually available for the
particular RG map used in a constructive argument?

\subsection{Categorical semantics and functorial QFT}

Institution theory formalizes invariance of truth under change of notation
\cite{GoguenBurstall1992}.  Locally covariant QFT represents a theory as a
functor from spacetimes to algebras
\cite{BrunettiFredenhagenVerch2003}; Fewster shows that functorial covariance
alone does not exhaust the criterion for ``the same physics'' across
spacetimes \cite{Fewster2015}.  These frameworks are mathematically richer
than our generic package in their intended domains.  Our contribution is not
a replacement category of physical theories.  It is a claim-centered overlay
for heterogeneous, often noninvertible proof interfaces.  Category language
is used only for the exact subcalculus and composable typed maps.

\subsection{Relational and informational interpretations}

Rovelli's relational quantum mechanics treats facts as relative to physical
systems and uses information-theoretic language in reconstructing quantum
mechanics \cite{Rovelli1996}.  Our relations are between descriptions,
observables, and proof obligations.  We remain neutral on whether quantum
states or events are observer-relative.  ``Information'' is eliminable
shorthand unless a bearer, equivalence relation, and inferential role are
specified.

\subsection{Novelty boundary}

No individual element is claimed without precedent: property-relative
equivalence, functorial maps, error moduli, reconstruction theorems, uniform
limits, and spectral positivity are established ideas.  The contribution is
their integration into a certificate profile designed for property loss and
recovery across constructive-QFT interfaces, together with a theorem that
changes the property's bearer and countermodels that force its central
hypotheses.  The literature search found no identical integrated framework,
but the paper makes no absolute priority claim.
